\centerline{\bf \Huge Полюсы и поляры }

\begin{flushright}
        \scriptsize{}
\end{flushright}

Рассмотрим окружность $\omega$ с центром $O$ и точку $P$, отличную от $O$. \textbf{Полярой} точки $P$ относительно $\omega$ называется прямая $\ell$, которая проходит через инверсную точку $P^{\prime}$ перпендикулярно прямой $O P$. Точка $P$ называется \textbf{полюсом} прямой $\ell$.

Везде далее мы будем считать, что мы находимся на проективной плоскости. Доопределим понятие поляры для $O$ и бесконечно удалённых точек. Полярой $O$ будем считать бесконечно удалённую прямую. Полярой бесконечно удалённой точки будем считать прямую, которая проходит через центр окружности перпендикулярно направлению, соответствующему бесконечно удалённой точке. Несложно видеть, что получилась биекция между точками и прямыми на проективной плоскости.

1. (a) \textbf{Основное свойство поляры.} Покажите, что точка $A$ лежит на поляре точки $B$ тогда и только тогда, когда $B$ лежит на поляре точки $A$.

\quad(b) Поляры точек $A$ и $B$ пересекаются в точке $C$. Чем является поляра точки $C$?

2. Через произвольную точку $X$ провели секущую к окружности $\omega$. Она пересекла окружность в точках $A$ и $B$. Докажите, что точка пересечения касательных к $\omega$ в точках $A$ и $B$ лежит на поляре $X$.

3. Докажите, что точка $A$ лежит на поляре точки $B$ относительно окружности с центром $O$ и радиусом $R$ тогда и только тогда, когда $(\overline{O A}, \overline{O B})=R^2$. Здесь в скобках написано скалярное произведение векторов.

4. Докажите, что точки $A, B$ и $C$ лежат на одной прямой тогда и только тогда, когда их поляры $a, b$ и $c$ пересекаются в одной точке.

5. (a) Пусть точка $Y$ лежит на поляре к точке $X$ относительно окружности $\omega$. Пусть прямая $X Y$ пересекает $\omega$ в точках $A$ и $B$. Докажите, что $(X, Y, A, B)=-1$.

\quad(b) \textbf{Полярное свойство секущих.} Через точку $X$ провели две прямые, которые пересекли окружность $\omega$ в точках $A, B$ и $C, D$ соответственно. Тогда точка пересечения $A C$ и $B D$ и точка пересечения $A D$ и $B C$ лежат на поляре точки $X$.